\documentclass[11pt]{article}

% ============================================================================
% PACKAGES
% ============================================================================
\usepackage[margin=1in]{geometry}
\usepackage{amsmath,amssymb,amsthm}
\usepackage{graphicx}
\usepackage{hyperref}
\usepackage{booktabs}
\usepackage{siunitx}
\usepackage[version=4]{mhchem}
\usepackage{fancyhdr}
\usepackage{enumitem}
\usepackage{mathtools}
\usepackage{braket}
\usepackage{xcolor}
\usepackage{framed}
\usepackage{tcolorbox}
\usepackage{physics}
\usepackage{float}

% ============================================================================
% PAGE SETUP
% ============================================================================
\pagestyle{fancy}
\fancyhf{}
\fancyhead[L]{\small CHM328 Companion Notes -- Lecture 2}
\fancyhead[R]{\small Alastair J.\ Price}
\fancyfoot[C]{\thepage}
\renewcommand{\headrulewidth}{0.4pt}

\hypersetup{
    colorlinks=true,
    linkcolor=blue!60!black,
    citecolor=blue!60!black,
    urlcolor=blue!60!black,
}

% Number equations within sections
\numberwithin{equation}{section}

\title{\textbf{Density Functional Theory:\\Transition States, Kinetics, and Selectivity}\\[1em]
\large Companion Notes for CHM328H1\,S -- Modern Physical Chemistry\\
University of Toronto, Winter 2026}
\author{Alastair J.\ Price\\
\textit{Guest Lecturer}\\
\small Department of Chemistry, University of Toronto}
\date{March 2026}

\begin{document}
\maketitle
\thispagestyle{fancy}

\begin{abstract}
\noindent
These notes accompany the second of two guest lectures on Density Functional Theory for Professor von Lilienfeld's CHM328H1\,S course. Building on the electronic structure and statistical mechanical foundations from Lecture~1, this lecture addresses the computational study of chemical reactions. Topics include the characterization of transition states as first-order saddle points on the potential energy surface, practical methods for locating and verifying transition states, the Eyring equation connecting barrier heights to rate constants, the role of entropy in determining activation free energies, and the prediction of regio- and stereoselectivity from computed energy differences. Several case studies from organic chemistry illustrate the methodology, and a detailed method selection guide is provided for practical applications.
\end{abstract}

\tableofcontents
\newpage

% ============================================================================
\section{Chemical Reactions on the Potential Energy Surface}
\label{sec:PES_reactions}
% ============================================================================

A chemical reaction corresponds to a path on the potential energy surface (PES) connecting one minimum (the reactant) to another (the product). Along this path, the system passes through a high-energy region; the highest point on the lowest-energy path is the \emph{transition state} (TS). The energy difference between the TS and the reactant is the \emph{activation energy}, and it determines how fast the reaction proceeds.

\subsection{The Minimum Energy Path and Reaction Coordinate}

The minimum energy path (MEP) is the steepest-descent path from the transition state to the reactant and product minima in mass-weighted coordinates. The arc length along this path defines the \emph{intrinsic reaction coordinate} (IRC). A one-dimensional plot of energy versus the IRC produces the familiar reaction energy profile, with reactants in one well, products in another, and the transition state at the barrier maximum.

For a system with $M$ atoms, the PES has $3M-6$ dimensions. Even for a molecule as small as water ($M = 3$), the PES is three-dimensional, and visualization becomes impractical for anything larger. The computational challenge is to locate a single saddle point---the transition state---on this high-dimensional surface.

\subsection{Multi-Step Reactions}

Many reactions proceed through intermediates, with a separate transition state for each elementary step. The rate-determining step (RDS) is the one with the highest barrier measured from the lowest preceding intermediate (or from the resting state of the catalyst in catalytic cycles). When modelling a multi-step reaction, every minimum and every transition state must be separately optimized and characterized by a frequency calculation.


% ============================================================================
\section{Transition State Theory}
\label{sec:TST}
% ============================================================================

\subsection{Definition of a Transition State}

A transition state is defined rigorously as a first-order saddle point on the PES. This is a stationary point ($\nabla E = \mathbf{0}$) where the Hessian matrix has exactly one negative eigenvalue. In terms of vibrational frequencies:

\begin{center}
\begin{tabular}{lcc}
\toprule
& \textbf{Minimum} & \textbf{Transition State} \\
\midrule
Gradient & $\mathbf{0}$ & $\mathbf{0}$ \\
Hessian eigenvalues & All positive & Exactly one negative \\
Imaginary frequencies & 0 & 1 \\
\bottomrule
\end{tabular}
\end{center}

The eigenvector corresponding to the negative eigenvalue describes the atomic displacements along the reaction coordinate---the motion that converts reactants into products. For a bond-breaking reaction, this mode shows the two atoms moving apart; for a hydrogen transfer, it shows the hydrogen migrating from donor to acceptor.

If two imaginary frequencies are found, the structure is a second-order saddle point and is not a physically meaningful transition state. If zero imaginary frequencies are found, the structure is a minimum, not a TS.

\subsection{Assumptions of Transition State Theory}

Transition state theory (TST), developed by Eyring, Polanyi, and Evans in 1935, rests on five assumptions:

\begin{enumerate}
\item The Born--Oppenheimer approximation is valid.
\item Nuclear motion along the reaction coordinate is classical (no tunnelling).
\item A quasi-equilibrium exists between the reactants and the transition state: the TS is treated as a thermodynamic species in equilibrium with the reactants.
\item The transition state is crossed only once (no recrossing). TST therefore provides an upper bound to the true rate constant.
\item The reaction coordinate at the TS can be separated from the other $3M-7$ vibrational degrees of freedom.
\end{enumerate}

\subsection{The Eyring Equation}
\label{sec:eyring}

Under these assumptions, the rate constant is given by the Eyring equation:
\begin{equation}
\boxed{k(T) = \kappa\,\frac{k_BT}{h}\,e^{-\Delta G^\ddagger/RT},}
\label{eq:eyring}
\end{equation}
where:
\begin{itemize}
\item $\kappa$ is the transmission coefficient ($\leq 1$; often assumed equal to 1),
\item $k_BT/h \approx \SI{6.2e12}{s^{-1}}$ at \SI{298}{K} is the universal frequency factor,
\item $\Delta G^\ddagger = G(\text{TS}) - G(\text{reactants})$ is the Gibbs free energy of activation.
\end{itemize}

An equivalent form that separates enthalpic and entropic contributions is
\begin{equation}
k(T) = \kappa\,\frac{k_BT}{h}\,e^{\Delta S^\ddagger/R}\,e^{-\Delta H^\ddagger/RT}.
\label{eq:eyring_HS}
\end{equation}
Comparison with the Arrhenius equation $k = Ae^{-E_a/RT}$ yields $E_a \approx \Delta H^\ddagger + RT$.

\subsection{Derivation Sketch}

The derivation proceeds from the quasi-equilibrium assumption. The concentration of transition state species is related to the reactant concentration by
\begin{equation}
K^\ddagger = \frac{[\text{TS}]}{[\text{R}]} = \frac{q^\ddagger}{q_\text{R}}\,e^{-\Delta E_0^\ddagger/k_BT},
\end{equation}
where $q^\ddagger$ and $q_\text{R}$ are molecular partition functions and $\Delta E_0^\ddagger$ is the zero-point-corrected electronic energy difference.

The TS partition function is factored into a one-dimensional translational component along the reaction coordinate (treated as a particle in a box of infinitesimal width $\delta$) and the remaining $3M-7$ degrees of freedom:
\begin{equation}
q^\ddagger = q_\text{rc}^\ddagger \times q_\text{rem}^\ddagger, \qquad q_\text{rc}^\ddagger = \frac{\sqrt{2\pi\mu k_BT}}{h}\,\delta.
\end{equation}
The rate of crossing the dividing surface is $\bar{v}/\delta = \sqrt{k_BT/2\pi\mu}/\delta$. Combining these factors, the rate constant becomes
\begin{equation}
k = \frac{k_BT}{h}\,\frac{q_\text{rem}^\ddagger}{q_\text{R}}\,e^{-\Delta E_0^\ddagger/k_BT} = \frac{k_BT}{h}\,e^{-\Delta G^\ddagger/RT},
\end{equation}
where the second equality follows from the standard thermodynamic relation between the equilibrium constant and $\Delta G^\ddagger$.

\subsection{Sensitivity to Barrier Height}

The exponential dependence in Eq.~\eqref{eq:eyring} makes reaction rates extraordinarily sensitive to barrier heights. Table~\ref{tab:barriers} illustrates this at \SI{298}{K}.

\begin{table}[h]
\centering
\caption{Rate constants and half-lives as a function of barrier height at \SI{298}{K}.}
\label{tab:barriers}
\begin{tabular}{S[table-format=2.0]S[table-format=1.1e2]l}
\toprule
{$\Delta G^\ddagger$ (kcal/mol)} & {$k$ (s$^{-1}$)} & {Half-life} \\
\midrule
5 & 1.4e9 & ${\sim}\SI{0.5}{ns}$ \\
10 & 3.1e5 & ${\sim}\SI{2}{\micro s}$ \\
15 & 6.8e1 & ${\sim}\SI{10}{ms}$ \\
20 & 1.5e-2 & ${\sim}\SI{46}{s}$ \\
25 & 3.3e-6 & ${\sim}\SI{2.4}{days}$ \\
30 & 7.3e-10 & ${\sim}\SI{30}{years}$ \\
35 & 1.6e-13 & ${\sim}\SI{140000}{years}$ \\
\bottomrule
\end{tabular}
\end{table}

A \SI{1}{kcal/mol} error in $\Delta G^\ddagger$ changes the predicted rate by a factor of ${\sim}5.5$; a \SI{2}{kcal/mol} error changes it by a factor of ${\sim}30$. Chemical accuracy ($\SI{1}{kcal/mol}$) is therefore the minimum standard for meaningful rate predictions.


% ============================================================================
\section{Finding Transition States}
\label{sec:finding_TS}
% ============================================================================

Locating transition states is substantially more difficult than locating minima. A minimum is a ``bowl'' on the PES---any reasonable starting geometry will roll downhill to the nearest one. A transition state is a saddle---it requires simultaneously moving uphill along the reaction coordinate and downhill in all other directions.

\subsection{Method 1: Relaxed PES Scan}

The simplest approach is to choose a geometric coordinate that changes during the reaction (typically a breaking or forming bond length), fix it at a series of values between the reactant and product values, and optimize all remaining coordinates at each point. The maximum of the resulting energy profile provides an approximate TS geometry that can serve as a starting point for a proper TS optimization.

This method is intuitive and provides a visual picture of the reaction path, but it relies on identifying the correct scan coordinate. For reactions where multiple bonds change simultaneously, a one-dimensional scan may miss the true reaction path.

\subsection{Method 2: Synchronous Transit (QST2/QST3)}

The quadratic synchronous transit method interpolates a path between the reactant and product geometries, finds the maximum along this path, and refines it using eigenvector following.

\begin{description}
\item[QST2] requires only the reactant and product geometries. It constructs a quadratic path and locates its maximum as an initial TS guess.
\item[QST3] additionally accepts an initial TS guess (e.g., from a relaxed scan), which generally improves convergence.
\end{description}

A critical practical requirement: the atom ordering must be identical in the reactant, TS guess, and product input structures. The algorithm maps atoms by their index; inconsistent ordering produces nonsensical interpolations.

In Gaussian: \texttt{opt=(QST2)} or \texttt{opt=(QST3)}.

\subsection{Method 3: Nudged Elastic Band (NEB)}

The nudged elastic band method optimizes an entire chain of structures (``images'') connecting the reactant and product, connected by virtual springs. The optimization minimizes forces perpendicular to the path while the springs maintain approximately even spacing along it. The climbing-image variant (CI-NEB) drives the highest-energy image uphill along the path to converge on the true saddle point.

NEB does not require a TS guess, finds the minimum energy path rather than just the TS, and handles complex reactions where the reaction coordinate is not obvious. Its disadvantage is the computational cost of optimizing many images simultaneously. In ORCA: \texttt{NEB-TS}.

\subsection{Method 4: Eigenvector Following (Berny Algorithm)}

The most direct approach: starting from a geometry near the TS, compute the Hessian, identify the eigenvector with the most negative eigenvalue, maximize along that direction, and minimize along all others. The algorithm iterates until convergence.

In Gaussian: \texttt{opt=(TS,CalcFC)} (Hessian at first step) or \texttt{opt=(TS,CalcAll)} (Hessian at every step). \texttt{CalcAll} is more robust but more expensive.

The critical requirement is that the starting geometry must already be close to the TS. If it is not, the optimizer may converge to a different saddle point or fail entirely.

\subsection{Summary of Methods}

\begin{table}[h]
\centering
\caption{Comparison of transition state search methods.}
\label{tab:TS_methods}
\begin{tabular}{llll}
\toprule
\textbf{Method} & \textbf{Input Required} & \textbf{Strengths} & \textbf{Weaknesses} \\
\midrule
Relaxed scan & Scan coordinate & Intuitive, visual & 1D only \\
QST2/QST3 & R \& P (+ TS guess) & Automated & Atom ordering \\
NEB/CI-NEB & R \& P & No TS guess needed & Expensive \\
Eigenvector following & Near-TS geometry & Fast convergence & Needs good guess \\
\bottomrule
\end{tabular}
\end{table}


% ============================================================================
\section{Verification of Transition States}
\label{sec:verification}
% ============================================================================

After locating a candidate transition state, verification is mandatory. This cannot be overemphasized: skipping verification is the single most common source of errors in published computational reaction studies.

\subsection{Step 1: Frequency Calculation}

Run a frequency calculation on the optimized TS structure. The result must show \emph{exactly one imaginary frequency}. Zero imaginary frequencies indicate a minimum; two or more indicate a higher-order saddle point.

\subsection{Step 2: Visualize the Imaginary Mode}

The eigenvector of the imaginary frequency shows the atomic motion along the reaction coordinate. This motion must correspond to the bond-breaking and bond-forming events expected for the reaction under study. If the imaginary mode instead shows, for example, a methyl group rotation or an out-of-plane bending, the wrong saddle point has been found.

\subsection{Step 3: Intrinsic Reaction Coordinate (IRC)}

The IRC follows the steepest-descent path from the TS in both directions (toward reactants and toward products):
\begin{equation}
\frac{d\mathbf{q}}{ds} = -\frac{\nabla V(\mathbf{q})}{|\nabla V(\mathbf{q})|},
\label{eq:IRC}
\end{equation}
where $s$ is the arc length and $\mathbf{q}$ are mass-weighted coordinates.

The forward IRC must reach the products; the backward IRC must reach the reactants. This proves that the TS connects the correct endpoints. Without an IRC, one does not know what reaction the TS describes.

In Gaussian: \texttt{IRC=(CalcFC,MaxPoints=50)}. In ORCA: the \texttt{IRC} block.

There are published examples---even in respected journals---where a transition state was reported for a claimed reaction, but the IRC actually connected to different species. The cost of running an IRC (a few minutes of computer time) is trivial compared to the cost of publishing incorrect results.


% ============================================================================
\section{Activation Parameters}
\label{sec:activation}
% ============================================================================

The activation parameters are obtained from DFT frequency calculations on both the reactants and the transition state, at the same level of theory:
\begin{align}
\Delta E^\ddagger &= E_\text{elec}(\text{TS}) - E_\text{elec}(\text{R}), \label{eq:dEddagger}\\
\Delta E_0^\ddagger &= \Delta E^\ddagger + \Delta E_\text{ZPE}^\ddagger, \label{eq:dE0ddagger}\\
\Delta H^\ddagger &= \Delta E_0^\ddagger + \Delta H_\text{thermal}^\ddagger, \label{eq:dHddagger}\\
\Delta G^\ddagger &= \Delta H^\ddagger - T\Delta S^\ddagger. \label{eq:dGddagger}
\end{align}

The transition state has $3M-7$ real vibrational modes (the imaginary mode, which is the reaction coordinate, is excluded from the vibrational partition function). This is handled automatically by quantum chemistry software.

The distinction between these quantities matters:
\begin{itemize}
\item $\Delta E^\ddagger$ is the raw electronic energy barrier.
\item $\Delta E_0^\ddagger$ includes zero-point energy corrections (typically lowering the barrier by \SIrange{1}{3}{kcal/mol}).
\item $\Delta H^\ddagger$ adds thermal corrections.
\item $\Delta G^\ddagger$ includes entropy and is the quantity that enters the Eyring equation.
\end{itemize}

Only $\Delta G^\ddagger$ should be used for computing rate constants.


% ============================================================================
\section{Entropy Effects on Activation Barriers}
\label{sec:entropy}
% ============================================================================

The entropic contribution to the activation free energy, $-T\Delta S^\ddagger$, is frequently the dominant correction that separates $\Delta G^\ddagger$ from $\Delta E^\ddagger$. Its sign and magnitude depend on the molecularity of the reaction.

\subsection{Classification by Molecularity}

\begin{table}[h]
\centering
\caption{Entropy effects on activation barriers classified by reaction type.}
\label{tab:entropy}
\begin{tabular}{lcccl}
\toprule
\textbf{Type} & $\Delta n^\ddagger$ & $\Delta S^\ddagger$ & \textbf{Effect on} $\Delta G^\ddagger$ & \textbf{Example} \\
\midrule
Unimolecular & 0 & $\approx 0$ & Small & Isomerization \\
Bimolecular (tight TS) & $-1$ & $-40$ to $-60$ J/(mol$\cdot$K) & Raises by 10--15 kcal/mol & \ce{S_N2}, Diels--Alder \\
Dissociative & $+1$ & $+40$ to $+60$ J/(mol$\cdot$K) & Lowers by 10--15 kcal/mol & Bond homolysis \\
\bottomrule
\end{tabular}
\end{table}

\subsection{Physical Origin}

When two molecules combine to form a single transition state complex, translational and rotational degrees of freedom are lost. Each lost translational degree of freedom costs approximately \SI{35}{J/(mol\cdot K)} of entropy, which at \SI{298}{K} corresponds to $T\Delta S \approx \SI{2.4}{kcal/mol}$. The total translational entropy penalty for a bimolecular association is approximately \SIrange{10}{15}{kcal/mol} at room temperature.

For dissociative reactions, the reverse occurs: the TS is looser than the reactant, and the gain in translational freedom lowers $\Delta G^\ddagger$ relative to $\Delta E^\ddagger$.

For unimolecular reactions (e.g., isomerizations), the number of molecules does not change between reactant and TS, so $\Delta S^\ddagger \approx 0$ and $\Delta G^\ddagger \approx \Delta H^\ddagger \approx \Delta E^\ddagger$ (roughly).

\subsection{Implications}

The free energy barrier for a bimolecular reaction is \emph{always} higher than the electronic energy barrier. This is a direct consequence of the entropic penalty of assembling two separate molecules into a single ordered transition state. Conversely, dissociative processes are entropically favoured, which explains why many bond-cleavage reactions become faster at high temperature despite having substantial electronic barriers.

This is one of the most important practical insights in computational reaction chemistry: \emph{the electronic energy alone is not predictive of reaction rates for bimolecular reactions.}


% ============================================================================
\section{Case Studies}
\label{sec:case_studies}
% ============================================================================

\subsection{Case Study 1: \texorpdfstring{S\textsubscript{N}2}{SN2} Reaction}

The prototypical bimolecular nucleophilic substitution:
\begin{equation}
\ce{Cl^- + CH3Br -> CH3Cl + Br^-}.
\end{equation}

\paragraph{Mechanism.} The nucleophile (\ce{Cl^-}) attacks the electrophilic carbon from the back side while the leaving group (\ce{Br^-}) departs. The carbon undergoes Walden inversion (the three hydrogens ``flip'' like an umbrella). The reaction is concerted---one transition state, no intermediate.

\paragraph{Transition state geometry.} The TS has a trigonal bipyramidal carbon with the nucleophile and leaving group occupying the axial positions, partially bonded to carbon. The \ce{C-Cl} bond is elongated relative to the product; the \ce{C-Br} bond is elongated relative to the reactant.

\paragraph{DFT Results} (B3LYP-D3(BJ)/def2-TZVP):
\begin{itemize}
\item Electronic barrier: $\Delta E^\ddagger \approx \SI{5}{kcal/mol}$ (gas phase).
\item ZPE-corrected barrier: $\Delta E_0^\ddagger \approx \SI{4.8}{kcal/mol}$.
\item Free energy barrier: $\Delta G^\ddagger \approx \SI{23}{kcal/mol}$ (gas phase, \SI{298}{K}).
\end{itemize}

\paragraph{Entropy analysis.} The enormous gap between $\Delta E^\ddagger$ (\SI{5}{kcal/mol}) and $\Delta G^\ddagger$ (\SI{23}{kcal/mol}) arises almost entirely from the entropic penalty of assembling two species (\ce{Cl^-} and \ce{CH3Br}) into a single TS complex. The translational entropy loss alone contributes approximately \SI{18}{kcal/mol} to $\Delta G^\ddagger$.

\paragraph{Eyring calculation.} At \SI{298}{K}:
\begin{align}
k &= \frac{k_BT}{h}\,e^{-\Delta G^\ddagger/RT} \nonumber\\
&= (\SI{6.2e12}{s^{-1}})\,e^{-96000/(8.314 \times 298)} \nonumber\\
&= (\SI{6.2e12}{s^{-1}})\times(\SI{1.4e-17}{}) \nonumber\\
&\approx \SI{8.4e-5}{s^{-1}}.
\end{align}
This predicts a slow gas-phase reaction. In solution, the polar solvent stabilizes the charged TS, lowering $\Delta G^\ddagger$ by \SIrange{8}{10}{kcal/mol} and increasing the rate by many orders of magnitude.

\paragraph{Solvent effects.} Gas-phase $\Delta G^\ddagger \approx \SI{23}{kcal/mol}$; aqueous $\Delta G^\ddagger \approx \SI{14}{kcal/mol}$. Using Eyring: $k_\text{aq} \approx \SI{e3}{s^{-1}}$, consistent with the experimental observation that \ce{S_N2} reactions are fast in polar solvents. This is one of the clearest illustrations of why implicit solvation (SMD or PCM) is essential for ionic reactions.

\subsection{Case Study 2: Diels--Alder Cycloaddition}

\begin{equation}
\ce{C4H6 + C2H4 -> C6H10} \quad (\text{butadiene + ethene} \to \text{cyclohexene}).
\end{equation}

\paragraph{Mechanism.} The Diels--Alder reaction is a concerted [4+2] pericyclic cycloaddition. Two new \ce{C-C} bonds form simultaneously in a cyclic transition state. The TS is aromatic in character: six electrons are delocalized in a cyclic array, satisfying the H\"{u}ckel $4n+2$ rule.

\paragraph{DFT Results} (B3LYP-D3/6-311+G(d,p)):
\begin{itemize}
\item $\Delta E^\ddagger \approx \SI{22}{kcal/mol}$
\item $\Delta H^\ddagger \approx \SI{20}{kcal/mol}$ (after ZPE correction)
\item $\Delta G^\ddagger \approx \SI{33}{kcal/mol}$
\item $\Delta G_\text{rxn} \approx \SI{-38}{kcal/mol}$ (strongly exergonic)
\end{itemize}

\paragraph{Entropy analysis.} The gap $\Delta G^\ddagger - \Delta E^\ddagger \approx \SI{11}{kcal/mol}$ arises from the exceptionally negative activation entropy: $\Delta S^\ddagger \approx \SI{-45}{J/(mol\cdot K)}$. Two molecules combine into a highly ordered cyclic TS, losing both translational and rotational freedom. At \SI{298}{K}:
\begin{equation}
-T\Delta S^\ddagger = -(298)(-0.045) \approx \SI{+13.4}{kcal/mol}.
\end{equation}
This entropic penalty explains why Diels--Alder reactions require elevated temperatures despite being strongly exergonic---the thermodynamic driving force is large, but the kinetic barrier (dominated by entropy) is substantial.

\paragraph{Comparison with experiment.} The experimental activation energy for butadiene + ethene is approximately \SI{27}{kcal/mol}. The B3LYP value of $\sim\!\SI{24}{kcal/mol}$ for $\Delta H^\ddagger$ is in reasonable agreement; higher-rung functionals such as $\omega$B97X-D or M06-2X give improved barriers.

\subsection{Case Study 3: Electrophilic Aromatic Substitution (Toluene Nitration)}

The nitration of toluene illustrates regioselectivity prediction:
\begin{equation}
\ce{C6H5CH3 + NO2+ -> products}.
\end{equation}

The methyl group is an ortho/para director. DFT can compute the activation barriers for \ce{NO2+} attack at each position:

\begin{table}[H]
\centering
\caption{Computed barriers and product distribution for toluene nitration.}
\label{tab:toluene}
\begin{tabular}{lccc}
\toprule
\textbf{Position} & $\Delta G^\ddagger$ (kcal/mol) & \textbf{Relative rate} & \textbf{Product \%} \\
\midrule
\textit{ortho} & 21.5 & 1.0 & ${\sim}58\%$ \\
\textit{meta} & 23.8 & 0.02 & ${\sim}4\%$ \\
\textit{para} & 21.8 & 0.6 & ${\sim}37\%$ \\
\bottomrule
\end{tabular}
\end{table}

The statistical factor matters: there are two equivalent \textit{ortho} positions but only one \textit{para} position. Even though the \textit{para} barrier is marginally lower than the individual \textit{ortho} barrier, the factor-of-two statistical advantage gives \textit{ortho} a larger share of the product distribution.

The computed product ratios (${\sim}58:4:37$) are in good agreement with experimental observations, demonstrating that DFT can quantitatively reproduce substituent effects on regioselectivity.


% ============================================================================
\section{Selectivity}
\label{sec:selectivity}
% ============================================================================

When a reaction can produce multiple products through different transition states, selectivity is determined by the difference in barrier heights, $\Delta\Delta G^\ddagger$.

\subsection{The Selectivity Equation}

The ratio of rate constants (and hence product ratios under kinetic control) is
\begin{equation}
\boxed{\frac{k_A}{k_B} = e^{-\Delta\Delta G^\ddagger/RT},}
\label{eq:selectivity}
\end{equation}
where $\Delta\Delta G^\ddagger = \Delta G^\ddagger_A - \Delta G^\ddagger_B$ is the difference in activation free energies. The relationship between $\Delta\Delta G^\ddagger$ and selectivity at \SI{298}{K} is shown in Table~\ref{tab:selectivity}.

\begin{table}[h]
\centering
\caption{Relationship between barrier difference and selectivity at \SI{298}{K}.}
\label{tab:selectivity}
\begin{tabular}{S[table-format=1.1]l}
\toprule
{$\Delta\Delta G^\ddagger$ (kcal/mol)} & {Product ratio (A\,:\,B)} \\
\midrule
0.5 & 70\,:\,30 \\
1.0 & 84\,:\,16 \\
1.4 & 90\,:\,10 \\
2.0 & 97\,:\,3 \\
3.0 & 99.4\,:\,0.6 \\
\bottomrule
\end{tabular}
\end{table}

Exquisite selectivity ($>95\%$) corresponds to energy differences of only \SIrange{2}{3}{kcal/mol}---right at the edge of DFT accuracy. This makes selectivity prediction both powerful (the numbers are physically meaningful) and challenging (the error bars are comparable to the quantity of interest).

\subsection{Kinetic vs.\ Thermodynamic Control}

Under \emph{kinetic control} (low temperature, short reaction time, irreversible conditions), the product that forms fastest dominates---this is determined by the barrier heights $\Delta G^\ddagger$. Under \emph{thermodynamic control} (high temperature, long reaction time, reversible conditions), the most stable product dominates---this is determined by the reaction free energies $\Delta G_\text{rxn}$.

These can give different products. Consider two pathways: product A forms faster (lower $\Delta G^\ddagger_A$) but product B is more stable (lower $G_B$). At low temperature, A dominates (kinetic product); at high temperature, B dominates (thermodynamic product). The crossover temperature is where the product distribution reverses.

DFT predicts both the kinetic and thermodynamic products by computing both the barriers and the product energies.

\subsection{Enantioselectivity}

For asymmetric catalysis, the enantiomeric excess (ee) is:
\begin{equation}
\text{ee} = \frac{|k_R - k_S|}{k_R + k_S}\times 100\% = \tanh\!\left(\frac{\Delta\Delta G^\ddagger_{R/S}}{2RT}\right)\times 100\%.
\end{equation}
At \SI{298}{K}, $\Delta\Delta G^\ddagger = \SI{2}{kcal/mol}$ gives ${>}90\%$ ee, and \SI{3}{kcal/mol} gives ${>}99\%$ ee. These energy differences are within the reach of careful DFT calculations, making enantioselectivity prediction feasible---though it requires high-quality methodology and ideally benchmarking against experimental data.


% ============================================================================
\section{Method Selection Guide}
\label{sec:methods}
% ============================================================================

The choice of functional and basis set affects the accuracy of computed barriers and must be made deliberately. Table~\ref{tab:methods} provides a practical guide.

\begin{table}[h]
\centering
\caption{Method selection guide for computational reaction studies.}
\label{tab:methods}
\begin{tabular}{p{4.5cm}p{5.5cm}c}
\toprule
\textbf{Application} & \textbf{Recommended Method} & \textbf{Error} \\
\midrule
Quick screening, qualitative trends & B3LYP/6-31G(d) or GFN2-xTB & 3--5 kcal/mol \\
Standard organic reactions & B3LYP-D3(BJ)/def2-TZVP & 2--3 kcal/mol \\
General thermochemistry & $\omega$B97X-D/def2-TZVP & 1--3 kcal/mol \\
Transition metals, organometallics & PBE0-D3(BJ)/def2-TZVP & 2--4 kcal/mol \\
Charge-transfer, long-range & $\omega$B97X-D/def2-TZVP or CAM-B3LYP-D3 & 1--2 kcal/mol \\
High accuracy, enantioselectivity & DFT geometry + DLPNO-CCSD(T)/cc-pVTZ SP & ${<}1$ kcal/mol \\
Large molecules (${>}50$ atoms) & B3LYP-D3/6-31G(d) geom, def2-TZVP SP & 2--3 kcal/mol \\
\bottomrule
\end{tabular}
\end{table}

\subsection{Dispersion Corrections}

Dispersion corrections must be included for essentially all reaction studies involving molecules larger than a few atoms. The penalty for omitting dispersion is typically \SIrange{2}{5}{kcal/mol} in barrier heights for systems with significant non-covalent contacts. In Gaussian: \texttt{EmpiricalDispersion=GD3BJ}. In ORCA: \texttt{D3BJ} or \texttt{D4}. The computational overhead is negligible.

In my own research on the XDM dispersion model~\cite{PriceXDM2023,PriceDispersion2021}, I have found that the choice of damping function and the treatment of higher-order ($C_8$, $C_{10}$) dispersion coefficients can significantly affect the accuracy of computed interaction energies, particularly for systems with both covalent and non-covalent interactions simultaneously. For routine applications, D3(BJ) is a reliable default.

\subsection{Solvent Effects}

For reactions in solution---particularly those involving charged species---implicit solvation is essential. The SMD model~\cite{SMD2009} (specifically parameterized for solvation free energies) is generally recommended over PCM for thermochemistry.

The magnitude of solvent effects on barriers:
\begin{itemize}
\item Ionic reactions: \SIrange{5}{15}{kcal/mol} change in $\Delta G^\ddagger$ between gas phase and polar solvent.
\item Reactions with significant charge redistribution: \SIrange{2}{5}{kcal/mol}.
\item Neutral reactions with minimal polarity change: ${<}\SI{2}{kcal/mol}$.
\end{itemize}

Neglecting solvation for an ionic reaction can produce errors of an order of magnitude or more in the predicted rate constant.

\subsection{The Quasi-RRHO Correction}

For reactions involving flexible molecules, low-frequency vibrational modes ($<\!\SI{100}{cm^{-1}}$) contribute unrealistically to the vibrational entropy under the harmonic approximation. Grimme's quasi-RRHO correction~\cite{GrimmeQRRHO2012} should be applied as standard practice. This is the default in ORCA; in Gaussian, post-processing with tools such as GoodVibes is recommended.


% ============================================================================
\section{Common Pitfalls}
\label{sec:pitfalls}
% ============================================================================

The following errors are encountered frequently in computational reaction studies. Each is avoidable with careful methodology.

\subsection{Pitfall 1: Not Verifying the Transition State}

A structure with one imaginary frequency is not necessarily the TS for the reaction of interest. Without an IRC, one cannot determine which reactant and product the TS connects. Published examples exist where a claimed TS for reaction A actually connected species involved in reaction B. Always perform: (i) frequency calculation $\to$ one imaginary frequency, (ii) visualization of the imaginary mode, and (iii) IRC in both directions.

\subsection{Pitfall 2: Comparing Energies at Different Levels of Theory}

Every species in a reaction energy profile must be computed with the same functional, basis set, dispersion correction, and solvation model. A B3LYP energy cannot be compared to a PBE energy; a gas-phase energy cannot be compared to an SMD energy. This seems obvious, but it is violated surprisingly often in practice, particularly in multi-step reactions where different species may have been computed at different times or by different people.

\subsection{Pitfall 3: Reporting $\Delta E$ Instead of $\Delta G$}

The electronic energy barrier $\Delta E^\ddagger$ does not enter the Eyring equation---$\Delta G^\ddagger$ does. For unimolecular reactions, $\Delta E^\ddagger \approx \Delta G^\ddagger$ is often a reasonable approximation. For bimolecular reactions, $\Delta G^\ddagger$ can exceed $\Delta E^\ddagger$ by \SIrange{10}{15}{kcal/mol} due to the entropic assembly penalty. Reporting only $\Delta E^\ddagger$ for a bimolecular reaction is misleading at best.

\subsection{Pitfall 4: Ignoring Conformational Sampling}

Flexible molecules have multiple conformers. If the lowest-energy conformer of the reactant is missed, the barrier will be underestimated (the reference energy is too high). If the lowest-energy TS conformer is missed, the barrier will be overestimated. For molecules with more than one rotatable bond, conformational searching (e.g., CREST, RDKit) is advisable.

\subsection{Pitfall 5: Neglecting Solvent Effects for Ionic Reactions}

As illustrated by the \ce{S_N2} case study, gas-phase and solution-phase barriers for ionic reactions can differ by ${>}\SI{10}{kcal/mol}$. For any reaction involving charged species, an implicit solvation model (SMD preferred) is mandatory.

\subsection{Pitfall 6: Low-Frequency Mode Errors}

Vibrational modes below $\sim\!\SI{100}{cm^{-1}}$ cause disproportionately large errors in $S_\text{vib}$ and hence $\Delta G$. Apply the quasi-RRHO correction (Grimme) or quasi-harmonic approximation (Truhlar) as standard practice.

\subsection{Pitfall 7: Neglecting Dispersion}

Omitting dispersion corrections systematically underestimates the stability of compact structures (including many transition states) relative to extended ones, introducing errors of several kcal/mol. There is no reason to omit dispersion in modern calculations; the cost is negligible.


% ============================================================================
\section{Beyond Classical Transition State Theory}
\label{sec:beyond_TST}
% ============================================================================

\subsection{Quantum Tunnelling}

Classical TST assumes that nuclei do not tunnel through the barrier. For light atoms---especially hydrogen---this assumption breaks down. Quantum tunnelling allows the reaction to proceed through the barrier rather than over it, making the rate faster than the classical TST prediction.

Tunnelling is most important for:
\begin{itemize}
\item Hydrogen transfer reactions,
\item Low temperatures,
\item Narrow, high barriers.
\end{itemize}

Experimentally, tunnelling manifests as anomalously large kinetic isotope effects ($k_H/k_D \gg 7$) and non-Arrhenius temperature dependence.

Simple one-dimensional corrections include:
\begin{description}
\item[Wigner correction:] $\kappa_W = 1 + \frac{1}{24}\left(\frac{h|\nu^\ddagger|}{k_BT}\right)^2$, which adds a tunnelling factor based on the imaginary frequency.
\item[Eckart barrier:] Fits a one-dimensional potential to the reactant, TS, and product energies and computes the transmission probability analytically.
\end{description}

For quantitative tunnelling rates, multidimensional approaches (small-curvature tunnelling, SCT) are required. A dramatic example: the hydrogen atom transfer in methylhydroxycarbene to form acetaldehyde occurs entirely by tunnelling below \SI{20}{K}---the classical rate at that temperature is effectively zero.

\subsection{Kinetic Isotope Effects}

Kinetic isotope effects (KIEs) are among the most powerful experimental tools for probing reaction mechanisms, and DFT can predict them straightforwardly.

A \emph{primary} KIE occurs when the bond to H/D is broken in the rate-determining step. The zero-point energy of a \ce{C-H} bond is higher than that of a \ce{C-D} bond (because $\nu \propto 1/\sqrt{m}$ and deuterium is heavier), so the ZPE correction to the barrier is smaller for the deuterated species. Classical TST predicts $k_H/k_D \approx 2$--$7$; values exceeding 7 suggest a tunnelling contribution.

A \emph{secondary} KIE ($k_H/k_D \approx 0.8$--$1.2$) occurs at positions not directly involved in bond breaking and reflects rehybridization at the transition state.

DFT predicts KIEs by running frequency calculations with the two isotopic masses and computing the difference in $\Delta G^\ddagger$. The Bigeleisen--Mayer equation provides the rigorous framework for these calculations.


% ============================================================================
\section{Complete Computational Workflow}
\label{sec:workflow}
% ============================================================================

The complete workflow for studying a reaction with DFT is:

\begin{enumerate}
\item \textbf{Optimize reactants.} Run geometry optimization + frequency calculation. Confirm zero imaginary frequencies (minimum). Read $G(\text{R})$.

\item \textbf{Optimize products.} Same procedure. Confirm minimum. Read $G(\text{P})$.

\item \textbf{Find and optimize the transition state.} Use a relaxed scan, QST2/3, or NEB to generate a TS guess, then refine with eigenvector following. Run frequency calculation.

\item \textbf{Verify the TS.} Confirm exactly one imaginary frequency. Visualize the imaginary mode. Run IRC in both directions. Confirm the TS connects the correct reactants and products.

\item \textbf{Optimize intermediates} (for multi-step reactions). Confirm each is a minimum.

\item \textbf{Compute thermodynamic and kinetic quantities.}
\begin{align}
\Delta G^\ddagger &= G(\text{TS}) - G(\text{R}), \\
\Delta G_\text{rxn} &= G(\text{P}) - G(\text{R}), \\
k(T) &= \frac{k_BT}{h}\,e^{-\Delta G^\ddagger/RT}.
\end{align}

\item \textbf{Optional: higher-level single points.} Optimize geometries at a lower level (e.g., B3LYP-D3/def2-TZVP) and compute single-point energies at a higher level (e.g., DLPNO-CCSD(T)/cc-pVTZ). Apply thermal corrections from the DFT frequencies.

\item \textbf{Build the energy profile.} Plot $\Delta G$ for all species along the reaction coordinate.
\end{enumerate}

All species must be computed at the same level of theory (same functional, basis set, dispersion correction, solvation model, temperature, and pressure).


% ============================================================================
\section*{References}
% ============================================================================
\addcontentsline{toc}{section}{References}

\begin{enumerate}[label={[\arabic*]},leftmargin=*]
\bibitem{SzaboOstlund} A.\ Szabo and N.\ S.\ Ostlund, \textit{Modern Quantum Chemistry: Introduction to Advanced Electronic Structure Theory} (Dover, 1996).

\bibitem{Jensen} F.\ Jensen, \textit{Introduction to Computational Chemistry}, 3rd ed.\ (Wiley, 2017).

\bibitem{Cramer} C.\ J.\ Cramer, \textit{Essentials of Computational Chemistry: Theories and Models}, 2nd ed.\ (Wiley, 2004).

\bibitem{KochHolthausen} W.\ Koch and M.\ C.\ Holthausen, \textit{A Chemist's Guide to Density Functional Theory}, 2nd ed.\ (Wiley-VCH, 2001).

\bibitem{ParrYang} R.\ G.\ Parr and W.\ Yang, \textit{Density-Functional Theory of Atoms and Molecules} (Oxford University Press, 1989).

\bibitem{Eyring1935} H.\ Eyring, J.\ Chem.\ Phys.\ \textbf{3}, 107 (1935).

\bibitem{GrimmeD3} S.\ Grimme, J.\ Antony, S.\ Ehrlich, and H.\ Krieg, J.\ Chem.\ Phys.\ \textbf{132}, 154104 (2010).

\bibitem{GrimmeD3BJ} S.\ Grimme, S.\ Ehrlich, and L.\ Goerigk, J.\ Comput.\ Chem.\ \textbf{32}, 1456 (2011).

\bibitem{GrimmeQRRHO2012} S.\ Grimme, Chem.\ Eur.\ J.\ \textbf{18}, 9955 (2012).

\bibitem{PriceDispersion2021} A.\ J.\ Price, K.\ R.\ Bryenton, and E.\ R.\ Johnson, J.\ Chem.\ Phys.\ \textbf{154}, 230902 (2021).

\bibitem{PriceXDM2023} A.\ J.\ Price, A.\ Otero-de-la-Roza, and E.\ R.\ Johnson, Chem.\ Sci.\ \textbf{14}, 1252 (2023).

\bibitem{SMD2009} A.\ V.\ Marenich, C.\ J.\ Cramer, and D.\ G.\ Truhlar, J.\ Phys.\ Chem.\ B \textbf{113}, 6378 (2009).

\bibitem{TruhlarQH} R.\ F.\ Ribeiro, A.\ V.\ Marenich, C.\ J.\ Cramer, and D.\ G.\ Truhlar, J.\ Phys.\ Chem.\ B \textbf{115}, 14556 (2011).

\bibitem{Fukui1981} K.\ Fukui, Acc.\ Chem.\ Res.\ \textbf{14}, 363 (1981). (IRC theory)

\bibitem{Gonzalez1990} C.\ Gonzalez and H.\ B.\ Schlegel, J.\ Phys.\ Chem.\ \textbf{94}, 5523 (1990). (IRC implementation)

\bibitem{Houk2003} K.\ N.\ Houk, P.\ H.-Y.\ Cheong, Nature \textbf{455}, 309 (2008). (DFT for organic reactions)

\bibitem{Schreiner2011} P.\ R.\ Schreiner, H.\ P.\ Reisenauer, D.\ Ley, D.\ Gerbber, C.-H.\ Wu, and W.\ D.\ Allen, Science \textbf{332}, 1300 (2011). (Tunnelling in methylhydroxycarbene)

\bibitem{BeckeJohnson2005} A.\ D.\ Becke and E.\ R.\ Johnson, J.\ Chem.\ Phys.\ \textbf{122}, 154104 (2005). (XDM model)

\bibitem{CCCBDB} NIST Computational Chemistry Comparison and Benchmark Database, \url{https://cccbdb.nist.gov/}.
\end{enumerate}

\end{document}
